% =============================================================================
% Statements of mock midterm 2 (Topics 1, 2A, 2B), written once and \input by
%   mock2-midterm-topics-1-2b.tex            (exam paper)
%   mock2-midterm-topics-1-2b-solutions.tex  (worked solutions)
% Every statement is a real exercise copied verbatim (rules/exercise-lists.md):
% source line first, then the full text. Allowed edits only: math in LaTeX, tables
% rebuilt, line breaks undone, currency written as \$.
% Needs: \srcline{...} defined by the including document.
% =============================================================================

% ---------------------------------------------------------------- Question 1.1
\newcommand{\QoneA}{%
\srcline{Higgins (Koski--Higgins--Mitton), \emph{Analysis for Financial Management} 13e, Ch.~1,
Exercise 3, p.~32 (PDF 53) $\cdot$ Answer in book: yes, pp.~397--398 (PDF 418--419)}

\noindent Answer true or false to each of the following. Briefly explain your reasoning for each answer.
\begin{enumerate}[label=\alph*.,leftmargin=2em]
  \item If a company gets into financial difficulty, it can use some of its shareholders' equity to pay
        its bills for a time.
  \item It is impossible for a firm to have a negative book value of equity without the firm going into
        bankruptcy.
  \item If a company increases its dividend, its net income will decrease.
  \item You can construct a sources and uses statement for 2020 if you have a company's balance sheets
        for year-end 2020 and 2021.
  \item The goodwill account on the balance sheet is an attempt by accountants to measure the benefits
        that result from a company's public relations efforts in the community.
  \item A reduction in an asset account is a source of cash, while a reduction in a liability account is
        a use of cash.
  \item A company's market value of equity must always be higher than its book value of equity.
\end{enumerate}}

% ---------------------------------------------------------------- Question 1.2
\newcommand{\QoneB}{%
\srcline{Higgins (Koski--Higgins--Mitton), \emph{Analysis for Financial Management} 13e, Ch.~8,
Exercise 1, p.~327 (PDF 348) $\cdot$ Answer in book: yes, pp.~421--422 (PDF 442--443)}

\noindent Is each of the following statements true or false? Explain your answers briefly.
\begin{enumerate}[label=\alph*.,leftmargin=2em]
  \item Using the same risk-adjusted discount rate to discount all future cash flows ignores the fact
        that the more distant cash flows are often riskier than cash flows occurring sooner.
  \item The cost of capital, or WACC, is not the correct discount rate to use for all projects
        undertaken by a firm.
  \item If you can borrow all of the money you need for a project at 6 percent, the cost of capital for
        this project is 6 percent.
  \item The best way to estimate the cost of debt capital for a firm is to divide the interest expense
        on the income statement by the interest-bearing debt on the balance sheet.
  \item One reliable estimate of a privately held firm's equity beta is the average of the equity betas
        of several publicly held competitors.
  \item Nominal cash flows should always be discounted using a nominal discount rate.
\end{enumerate}}

% ---------------------------------------------------------------- Question 2.1
\newcommand{\QtwoA}{%
\srcline{Berk, DeMarzo \& Harford, \emph{Fundamentals of Corporate Finance} 2e, Ch.~13, Exercises 13.1
and 13.3, p.~403 (PDF 434) $\cdot$ Answer in book: no}

\noindent\textbf{1.} MV Corporation has debt with market value of \$100 million, common equity with a
book value of \$100 million, and preferred stock worth \$20 million outstanding. Its common equity trades
at \$50 per share, and the firm has 6 million shares outstanding. What weights should MV Corporation use
in its WACC?

\smallskip\noindent\textbf{3.} Book Co. has 1 million shares of common equity with a par (book) value of
\$1, retained earnings of \$30 million, and its shares have a market value of \$50 per share. It also has
debt with a par value of \$20 million that is trading at 101\% of par.
\begin{enumerate}[label=\alph*.,leftmargin=2em]
  \item What is the market value of its equity?
  \item What is the market value of its debt?
  \item What weights should it use in computing its WACC?
\end{enumerate}}

% ---------------------------------------------------------------- Question 2.2
\newcommand{\QtwoB}{%
\srcline{Berk, DeMarzo \& Harford, \emph{Fundamentals of Corporate Finance} 2e, Ch.~13, Exercises
13.18--13.19, p.~405 (PDF 436) $\cdot$ Answer in book: no}

\noindent\textbf{18.} RiverRocks (whose WACC is 12\%) is considering an acquisition of Raft Adventures
(whose WACC is 15\%). What is the appropriate discount rate for RiverRocks to use to evaluate the
acquisition? Why?

\smallskip\noindent\textbf{19.} RiverRocks' purchase of Raft Adventures (from Problem 18) will cost \$100
million, but will generate cash flows that start at \$15 million in one year and then grow at 4\% per year
forever. What is the NPV of the acquisition?}

% ---------------------------------------------------------------- Question 2.3
\newcommand{\QtwoC}{%
\srcline{Berk, DeMarzo \& Harford, \emph{Fundamentals of Corporate Finance} 2e, Ch.~13, Exercise 13.15,
p.~404 (PDF 435) $\cdot$ Answer in book: no}

\noindent Growth Company's current share price is \$20 and it is expected to pay a \$1 dividend per share
next year. After that, the firm's dividends are expected to grow at a rate of 4\% per year.
\begin{enumerate}[label=\alph*.,leftmargin=2em]
  \item What is an estimate of Growth Company's cost of equity?
  \item Growth Company also has preferred stock outstanding that pays a \$2 per share fixed dividend. If
        this stock is currently priced at \$28, what is Growth Company's cost of preferred stock?
  \item Growth Company has existing debt issued three years ago with a coupon rate of 6\%. The firm just
        issued new debt at par with a coupon rate of 6.5\%. What is Growth Company's pre-tax cost of
        debt?
  \item Growth Company has 5 million common shares outstanding and 1 million preferred shares
        outstanding, and its equity has a total book value of \$50 million. Its liabilities have a
        market value of \$20 million. If Growth Company's common and preferred shares are priced as in
        parts (a) and (b), what is the market value of Growth Company's assets?
  \item Growth Company faces a 35\% tax rate. Given the information in parts (a)--(d), and your answers
        to those problems, what is Growth Company's WACC?
\end{enumerate}}

% ---------------------------------------------------------------- Question 2.4
\newcommand{\QtwoD}{%
\srcline{Berk, DeMarzo \& Harford, \emph{Fundamentals of Corporate Finance} 2e, Ch.~16, Exercise 16.3,
p.~491 (PDF 522); section instruction on p.~490 (PDF 521) $\cdot$ Answer in book: no}

\noindent\emph{Section instruction (Capital Structure in Perfect Capital Markets):} For problems in this
section assume no taxes or distress costs.

\smallskip\noindent Acort Industries owns assets that will have an 80\% probability of having a market
value of \$50 million in one year. There is a 20\% chance that the assets will be worth only \$20
million. The current risk-free rate is 5\%, and Acort's assets have a cost of capital of 10\%.
\begin{enumerate}[label=\alph*.,leftmargin=2em]
  \item If Acort is unlevered, what is the current market value of its equity?
  \item Suppose instead that Acort has debt with a face value of \$20 million due in one year. According
        to MM, what is the value of Acort's equity in this case?
  \item What is the expected return of Acort's equity without leverage? What is the expected return of
        Acort's equity with leverage?
  \item What is the lowest possible realized return of Acort's equity with and without leverage?
\end{enumerate}}

% ---------------------------------------------------------------- Question 3.1
\newcommand{\QthreeA}{%
\srcline{Berk, DeMarzo \& Harford, \emph{Fundamentals of Corporate Finance} 2e, Ch.~16, Exercise 16.10,
p.~492 (PDF 523) $\cdot$ Answer in book: no}

\noindent Suppose the corporate tax rate is 40\%. Consider a firm that earns \$1000 before interest and
taxes each year with no risk. The firm's capital expenditures equal its depreciation expenses each year,
and it will have no changes to its net working capital. The risk-free interest rate is 5\%.
\begin{enumerate}[label=\alph*.,leftmargin=2em]
  \item Suppose the firm has no debt and pays out its net income as a dividend each year. What is the
        value of the firm's equity?
  \item Suppose instead the firm makes interest payments of \$500 per year. What is the value of equity?
        What is the value of debt?
  \item What is the difference between the total value of the firm with leverage and without leverage?
  \item To what percentage of the value of the debt is the difference in part (c) equal?
\end{enumerate}}

% ---------------------------------------------------------------- Question 3.2
\newcommand{\QthreeB}{%
\srcline{Berk, DeMarzo \& Harford, \emph{Fundamentals of Corporate Finance} 2e, Ch.~16, Exercises
16.16--16.17, p.~493 (PDF 524) $\cdot$ Answer in book: no}

\noindent\textbf{16.} Milton Industries expects free cash flows of \$5 million each year. Milton's
corporate tax rate is 35\%, and its unlevered cost of capital is 15\%. The firm also has outstanding debt
of \$19.05 million, and it expects to maintain this level of debt permanently.
\begin{enumerate}[label=\alph*.,leftmargin=2em]
  \item What is the value of Milton Industries without leverage?
  \item What is the value of Milton Industries with leverage?
\end{enumerate}
\noindent\textbf{17.} NatNah, a builder of acoustic accessories, has no debt and an equity cost of
capital of 15\%. NatNah decides to increase its leverage to maintain a market debt-to-value ratio of 0.5.
Suppose its debt cost of capital is 9\% and its corporate tax rate is 35\%. If NatNah's pretax WACC
remains constant, what will its (effective after-tax) WACC be with the increase in leverage?}

% ---------------------------------------------------------------- Question 3.3
\newcommand{\QthreeC}{%
\srcline{Berk, DeMarzo \& Harford, \emph{Fundamentals of Corporate Finance} 2e, Ch.~16, Exercise 16.23,
p.~494 (PDF 525) $\cdot$ Answer in book: no}

\noindent Marpor Industries has no debt and expects to generate free cash flows of \$16 million each year.
Marpor believes that if it permanently increases its level of debt to \$40 million, the risk of financial
distress may cause it to lose some customers and receive less favorable terms from its suppliers. As a
result, Marpor's expected free cash flows with debt will be only \$15 million per year. Suppose Marpor's
tax rate is 35\%, the risk-free rate is 5\%, the expected return of the market is 15\%, and the beta of
Marpor's free cash flows is 1.10 (with or without leverage).
\begin{enumerate}[label=\alph*.,leftmargin=2em]
  \item Estimate Marpor's value without leverage.
  \item Estimate Marpor's value with the new leverage.
\end{enumerate}}

% ---------------------------------------------------------------- Question 3.4
\newcommand{\QthreeD}{%
\srcline{Berk, DeMarzo \& Harford, \emph{Fundamentals of Corporate Finance} 2e, Ch.~17, Exercise 17.9,
p.~528 (PDF 559) $\cdot$ Answer in book: no}

\noindent The HNH Corporation will pay a constant dividend of \$2 per share, per year, in perpetuity.
Assume all investors pay a 20\% tax on dividends and that there is no capital gains tax. The cost of
capital for investing in HNH stock is 12\%.
\begin{enumerate}[label=\alph*.,leftmargin=2em]
  \item What is the price of a share of HNH stock?
  \item Assume that management makes a surprise announcement that HNH will no longer pay dividends but
        will use the cash to repurchase stock instead. What is the price of a share of HNH stock now?
\end{enumerate}}
